\section{Data, Code \& Reproducibility}
\label{sec:availability}

This section documents the artifacts that constitute the Phase~B/C release: the raw capture corpus, the cryptographic chain logs, the rendered emission tiles, the derived analysis products, the trained verifier checkpoints, and the source code, together with the reproducibility entry points that tie them to the results reported above. The analysis code, the recording-protocol source with its third-party chain verifier, and this paper accompany the release in the PolieBotics GitHub organization (\url{https://github.com/poliebotics}); the patent filings are published in the companion repository \url{https://github.com/poliebotics/PolieBotics}. We also state explicitly that the corpus contains identifiable human-subject captures and therefore requires privacy-aware handling on any public distribution. The release is published with all rights reserved: no open-source license and no patent license is granted.

\subsection{Scope of the release}
\label{sec:availability:scope}

The release comprises two same-rig capture sessions of a single human performer, recorded under the Truth Beam protocol on one projector--camera rig \cite{patent2025}. Session~D2 (\texttt{cittadel\_first\_human}) contains \num{5992} chain rows; session~V10 (\texttt{cittadel\_v10\_mainnet\_improv}) contains \num{3743} chain rows. Each chain row $t$ corresponds to one physical exposure: a chain-derived emission $E_t$ projected onto the scene and captured as a raw BayerRG8 frame $C_t$ (\Cref{sec:protocol_chain}). Both sessions are complete in the sense that every chain row $t$ has both a rendered emission tile and a raw capture frame. D2 was recorded under the v9 chain and V10 under the v10 chain, so the release also exercises both chain-transition variants on a single rig.

We emphasize the scope of the corpus so that downstream users do not over-read it: this is a \emph{same-rig, two-session} dataset captured with one camera, one projector, and one performer. The generalization results reported in \Cref{sec:results_main_generalization} are same-rig cross-session (and same-rig cross-protocol); nothing in this release establishes cross-rig, cross-camera, cross-projector, or cross-subject behavior, and the dataset should not be described as supporting such claims.

\subsection{Raw capture frames}
\label{sec:availability:raw}

Raw capture frames are stored as headerless \texttt{uint8} BayerRG8 buffers, one file per chain row (\texttt{frame\_\{t:06d\}.raw}), where the file index $t$ is the chain row index rather than the camera's GenICam capture-frame counter. Each raw frame is exactly \num{24472000} bytes, corresponding to a single-channel $4600 \times 5320$ ($H \times W$) color-filter-array image ($4600 \times 5320 = \num{24472000}$ exactly), with the RGGB Bayer pattern documented in \Cref{sec:hardware}. The corpus is \SI{8}{bit} end-to-end.

The per-session raw totals follow directly from the frame counts and the exact per-frame size: D2 contributes $\num{5992} \times \num{24472000} = \num{146636224000}$ bytes ($\approx \SI{146.64}{\giga\byte}$), and V10 contributes $\num{3743} \times \num{24472000} = \num{91598696000}$ bytes ($\approx \SI{91.60}{\giga\byte}$), for a combined raw-Bayer footprint of \num{238234920000} bytes ($\approx \SI{238.23}{\giga\byte}$, \SI{221.90}{\gibi\byte}). The raw frames are read-only inputs and were never modified during analysis. For transfer, the raw Bayer compresses to roughly \SI{62}{\percent} of its original size with general-purpose codecs (a sample frame of \num{24472000} bytes compresses to \num{15074526} bytes with gzip and \num{15228672} bytes with zstd), giving a compressed transfer footprint near \SI{148}{\giga\byte}; the on-disk footprint after decompression is unchanged.

\subsection{Chain logs}
\label{sec:availability:chainlogs}

Each session ships one chain log, \texttt{chain\_log.csv}, recording the cryptographic chain state and per-row metadata defined in \Cref{sec:protocol_chain}. The D2 log uses the v9 thirteen-column schema (header columns $t$, \texttt{S\_t\_hex}, \texttt{bayer\_blake3\_hex}, \texttt{capture\_frame\_id}, \texttt{aravis\_device\_timestamp\_ns}, \texttt{capture\_wall\_ns}, \texttt{emission\_live\_pixel\_blake3\_hex}, \texttt{emission\_live\_queued\_wall\_ns}, \texttt{meta\_hex}, \texttt{emission\_png\_path}, \texttt{drand\_round\_number}, \texttt{drand\_round\_value\_hex}, \texttt{drand\_staleness\_ms}); the V10 log uses the v10 fifteen-column schema, which adds \texttt{ai\_payload\_count} and \texttt{ai\_payload\_root\_hex}. The logs are small: the D2 chain log is \num{2514982} bytes and the V10 chain log is \num{1822011} bytes, roughly \SI{4.3}{\mega\byte} combined. Despite their negligible size they are load-bearing for verification, since they carry the per-row chain state $S_t$, the raw-capture hash \texttt{bayer\_blake3\_hex}, the folded capture metadata, and the drand beacon fields required to rewalk the chain \cite{drandbeacon}.

\subsection{Rendered emission tiles}
\label{sec:availability:emissions}

The emissions actually projected during recording are released as rendered tiles at \texttt{\{session\}/derived/Emissions/tile\_\{t:06d\}.png}: $1080 \times 1920$ (PIL size $W{=}1920$, $H{=}1080$) RGB \SI{8}{bit}-per-channel PNGs, one per chain row. The set is complete for both sessions: \num{5992} tiles for D2 and \num{3743} for V10, a total of \num{9735} PNGs. The emission-tile set occupies \num{23816426030} bytes ($\approx \SI{23.82}{\giga\byte}$; D2 $\approx \SI{14.66}{\giga\byte}$, V10 $\approx \SI{9.16}{\giga\byte}$), with a mean of roughly \SI{2.45}{\mega\byte} per PNG. These tiles are the bit-exact targets $E_t$ used throughout the verifier and red-team experiments.

\subsection{Derived analysis products}
\label{sec:availability:derived}

To allow the per-frame and per-pixel results of \Cref{sec:results_main,sec:verifier_method} to be re-examined without re-running the verifier on the bulk corpus, the release includes a compact parsable data pack. A wide per-frame table, \texttt{visual\_metrics\_wide.csv} (\num{9735} rows $\times$ 33 columns, one row per chain frame across both sessions), records the per-frame red-area, robust-$z$, excess, delta, and availability statistics underlying the headline frame-level numbers; it is committed directly in the public code repository (\texttt{results/csv/}), alongside the raw EXP-6 per-frame ranking table (\texttt{results/eval/exp6\_correct\_e\_rank/}, $n=120$ frames $\times$ 51 candidates with per-candidate scores). Recomputing the frame-level contrasts from the compact table yields AUROC $0.9998$ for both (see the precision note in \Cref{sec:results_main_framelevel}). Alongside it, 208 compact NPZ map shards (D2: 120; V10: 88; $\approx 50$ frames each, $\approx \SI{14}{\giga\byte}$ total) carry per-frame robust-$z$ and excess maps and low-resolution RGB thumbnails, sufficient to redesign the visualizations of \Cref{sec:results_main} without touching the larger bulk shards.

These derived products are diagnostic and visualization artifacts; the verifier-tier quantitative claims rest on the Phase~G checkpoints and the protocol substrate, not on the compact maps. The NPZ shards contain low-resolution RGB thumbnails of the captures and therefore inherit the human-subject privacy considerations of \Cref{sec:availability:privacy}.

\subsection{Verifier checkpoints}
\label{sec:availability:checkpoints}

The trained Phase~G diffusion verifier (\Cref{sec:verifier_method}) is the load-bearing released model. Each checkpoint (\texttt{model\_final.pt}) is \num{477531127} bytes ($\approx \SI{478}{\mega\byte}$); per-step intermediate checkpoints are of comparable size (\num{477533091} bytes). The verifier-tier reference for the release visualizations is the \emph{main} model paired with its architecture- and data-matched \emph{shuffled} control, which together establish that the discriminative signal is the real chain coupling rather than generic emission statistics or an architecture artifact (\Cref{sec:verifier_method}). The Phase~G experiment directory as a whole, including the per-step training checkpoints across the \emph{main}, \emph{shuffled}, \emph{synthetic\_positive}, and \emph{lambda\_smoke} runs, is \SI{7.6}{\giga\byte}; the \emph{lambda\_smoke} run is a smoke test and is not a release model. We recommend that any public release deliberately select which checkpoints to distribute rather than ship the full working tree.

\subsection{Source code}
\label{sec:availability:code}

The analysis codebase under \texttt{src/} is small: \SI{1.3}{\mega\byte} on disk across 67 Python modules organized into 11 module directories (\texttt{models}, \texttt{data}, \texttt{phase\_f}, \texttt{eval}, \texttt{phase\_g}, \texttt{training}, \texttt{losses}, \texttt{preprocessing}, \texttt{phase\_h}, \texttt{utils}, \texttt{diagnostics}). The full working repository excluding the Python virtual environment, the data symlinks, and the experiment checkpoints is on the order of \SI{10}{\mega\byte} of code. The release upload bundle is sized as raw Bayer (\SI{238.23}{\giga\byte}) plus chain logs ($\approx \SI{4}{\mega\byte}$) plus emission tiles (\SI{23.82}{\giga\byte}) plus the code workspace ($\approx \SI{10}{\mega\byte}$), giving a total upload of approximately \SI{262.06}{\giga\byte}; the Python virtual environment (\SI{7.2}{\giga\byte}) is explicitly excluded. After a run, the intended retained download footprint is approximately \SI{3.5}{\giga\byte} (final-epoch checkpoints plus findings, results, logs, and TensorBoard summaries), with the per-experiment intermediate training checkpoints removed.

\subsection{Reproducibility entry points}
\label{sec:availability:repro}

The release is structured so that the central results can be reproduced from primary artifacts:

\begin{itemize}
  \item \textbf{Protocol and chain re-walk.} The chain logs (\Cref{sec:availability:chainlogs}) together with the per-channel BLAKE3-XOF seed derivation \cite{blake3} and the bit-exact emission renderer (\Cref{sec:protocol_chain}) regenerate each emission tile $E_t$ from the recorded chain state $S_t$. Stored emission tiles can thus be checked against a recomputed render, and the raw-capture hash \texttt{bayer\_blake3\_hex} can be recomputed from the released \texttt{frame\_\{t:06d\}.raw} buffers to confirm the closed feedback loop. Third-party chain-walk verifiers for both protocol versions ship with the code release (\texttt{verify\_v9.py} for D2, \texttt{verify\_v10.py} for V10); each recomputes $S_0$, every per-row transition, the post-final terminal state against the manifest's anchored $S_N$, and every pulse commitment, and each offers a \texttt{-{}-logs-only} mode that verifies the chain mathematics from the $\approx$\SI{4}{\mega\byte} of session metadata alone.
  \item \textbf{Verifier evaluation.} The Phase~G checkpoints (\Cref{sec:availability:checkpoints}), applied to released $(C_t, E_t)$ pairs at the fixed evaluation timesteps and noise count described in \Cref{sec:verifier_method}, regenerate the noise-prediction residual scores $\lVert \hat{\epsilon}-\epsilon\rVert^2$ from which the AUROC and $z$-statistics of \Cref{sec:results_main} are computed.
  \item \textbf{Visualization re-derivation.} The compact CSV and NPZ products (\Cref{sec:availability:derived}) reproduce the per-frame and per-pixel figures and the display recipe without re-running the model on the bulk corpus.
\end{itemize}

Two calibration items should be surfaced in any release manifest because they bound exact bit-level reproducibility downstream. First, the measured black level is unknown: the manifest records \texttt{black\_level: 0} with source \texttt{no\_measurement\_found}, and black-level subtraction is intentionally not applied (\Cref{sec:hardware}). Second, while the \emph{recording} convention is unambiguous---the emission committed at row $t$ is seeded from that same row's state, $\texttt{target\_chain\_row} = \texttt{capture\_row} = t$ (offset $0$, $E_t = \mathrm{gen}(\mathrm{XOF}(S_t))$; \Cref{subsec:pc_offset})---the empirically optimal frame-to-chain offset for the \emph{optical verifier} has not been separately measured on the held-out blocks, and auxiliary training and diagnostic pipelines have adopted shifted alignments; any offset-dependent reproduction should therefore run the offset diagnostic rather than assume a particular verifier-side alignment.

\subsection{Human-subject captures and privacy}
\label{sec:availability:privacy}

Both sessions feature an identifiable human performer: D2 is the first-human session (\texttt{cittadel\_first\_human}) and V10 is an improv session (\texttt{cittadel\_v10\_mainnet\_improv}). Consequently the raw frames, the rendered emission tiles, the derived RGB thumbnails in the NPZ shards, and any rendered previews all depict a person, and any \emph{public} release of these artifacts is a privacy and consent matter rather than a purely technical one. Consistent with the ethical-display rules applied to every promoted visualization in this work, no body or person region is manually suppressed, normalized away, or otherwise altered; the body-region treatment used in the analysis is an audit overlay only and never a display modification (\Cref{sec:results_redteam}). In the released sessions the performer is the author/operator himself, who consents to the publication of his own likeness. Public distribution of any other human-subject capture material would be gated on the performer's consent and handling, independently of the license terms of \Cref{sec:availability:license}.

\subsection{License}
\label{sec:availability:license}

This release is published with all rights reserved: no open-source license is granted for the source code, and no patent license is granted. As noted in \Cref{sec:availability:privacy}, these terms are distinct from and do not by themselves authorize redistribution of the human-subject capture material.
